\documentclass{article}

\usepackage{fullpage}
\usepackage{url}
\usepackage{graphicx}
\usepackage{amsmath}
\usepackage{physics}
\usepackage{pxfonts}
\usepackage{mathtools}

\DeclareMathOperator{\erfi}{erfi}
\newcommand{\R}{\ensuremath{\mathbb{R}}}

\title{Differential Geometry Excercise 1}
\author{Idan Stark}
\date{\today}

\begin{document}
\maketitle


\section*{Question 1}
Let $M$ and $N$ be smooth manifolds. For every chart $\varphi_i : U_i \rightarrow R^m$ where $U_i \subseteq M$ and chart $\varphi_j : U_j \rightarrow R^n$ where $U_j \subseteq N$, define a chart
\begin{equation}
    \varphi_{i,j}: U_i \times U_j \rightarrow R^{m+n}
    \qquad
    \varphi_{i,j}(p, q) = (\varphi_i(p), \varphi_j(q))
\end{equation}
where $ \left(\cdot,\cdot\right) $ denotes coordinates concatenation.
\newline
Now let $\{U_i, \varphi_i\}$ and $\{V_j, \psi_j\}$ be atlases in the respective structures of $M$ and $N$. Then the set
\begin{equation}
    A = \left\{\left(U_i \times V_j, \varphi_{i, j}\right)\right\}
\end{equation}
covers every point on $M \times N$. In addition, every chart in $A$ is smooth. So, for every two subsets of $M \times N$ that coincide, $(U_i \times V_j) \cap (U_k \times V_l) \ne \emptyset$, the transition map from one subset to the other
\begin{equation}
    \varphi_{k,l}\circ\varphi_{i,j}^{-1} = (\varphi_k\circ\varphi_i^{-1}, \varphi_l\circ\varphi_j^{-1})
\end{equation}
is also smooth. Thus, $A$ is an atlas and it generates a smooth structure over $M \times N$, making it a smooth manifold with dimension $\dim(M) + \dim(N)$.

\section*{Question 2}
Let $f: M \rightarrow N$ be an immersion that is also a submersion. Because it is an immersion, for every point $p \in M$ its differential $Df_p$ is injective. Because it is an submersion, for every point $p \in M$ its differential $Df_p$ is surjective. Thus, $Df_p$ is invertible for every point $p \in M$.
\newline
According to the Inverse Function Theorem, for a smooth function $f: U \rightarrow \R^n$ where $U \subset \R^n$ for which $Df_p$ is invertible at point $p \in U$, there exist neighborhoods $V, W$ of $p$ and $f(p)$ respectively such $f:V \rightarrow W$ is a local diffeomorphism.

\section*{Question 3}
Let $f = (f_1,..., f_k): M \rightarrow \R^k$ be a submersion. Therefore its differential $Df_p: T_p M \rightarrow \R^k$ at each point $p$ is surjective, and spans $\R^k$.
\newline
The differential in terms of coordinates $\left\{(df_i)_p: T_p M \rightarrow \R, i = 1,...,k\right\}$ are $k$ linear maps. Let us assume that they are linearly dependent. This means that there exist a constant vector $u \in \R^k, u \ne 0$ such that
\begin{equation}
    u \cdot \left(df_i\right)_p = 0
\end{equation}
So, for every vector $v \in T_p M$
\begin{equation}
    u \cdot \left(df_i\right)_p (v) = 0
\end{equation}
but since $Df_p$ is surjective
\begin{equation}
    \left\{\left(df_i\right)_p (v): v \in T_p M\right\} = \R^k
\end{equation}
and so, for every vector $w \in \R^k$, $u\cdot w = 0$, which means $u = 0$.
Therefore, the only linear combination of $(df_i)_p$ that vanished identically is the trivial one, and they are linearly independent.

\section*{Question 4}
$\R P^n$ is the set of all lines passing through the origin in $\R^{n+1}$:
\begin{equation}
    \R P^n = \left\{
        \{\lambda x: \lambda \ne 0\}
        : x \in \R^{n+1} \backslash \{0\}
    \right\}
\end{equation}
Then, let $U_i = \{\{\lambda x: \lambda \ne 0\}: x \in \R^{n+1} \backslash \{0\}, x_i \ne 0 \}$. For each such open set $U_i$, define a bijection $\varphi_i : U_i \rightarrow \R^{n}$ such that
\begin{equation}
    \varphi_i (\{\lambda x: \lambda \ne 0\}) = x_i^{-1} (x_1, ..., x_{i-1}, x_{i+1}, ..., x_{n+1})
\end{equation}
The inverse of this bijection is clear:
\begin{equation}
    \varphi_i^{-1} (x) = \{\lambda y : \lambda \ne 0\}, \qquad y = (x_1, ..., x_{i-1}, 1, x_{i+1}, ..., x_{n+1})
\end{equation}
The set of all charts $\{\varphi_i: i = 1,..., n+1 \}$ covers $\R P^n$. The image of each bijection $\varphi_i (U_i)$ is an open set, and $\varphi_i$ is a homeomorphism, therefore $\varphi_i(U_i \cap U_j) \quad \forall \quad i,j = 1,...,n+1$ is also an open set. Finally, for each $i, j = 1,...,n+1$:
\begin{equation}
    \varphi_i \circ \varphi_j^{-1} (x) = \begin{cases}
        x_i^{-1} (x_1, ..., x_{i-1}, x_{i+1}, ..., x_{j-1}, 1, x_{j+1},...,x_{n+1}) 
        & i < j
        \\
        x_i^{-1} (x_1, ..., x_{i-1}, x_{i+1},...,x_{n+1}) 
        & i = j
        \\
        x_i^{-1} (x_1, ..., x_{j-1}, 1, x_{j+1},..., x_{i-1}, x_{i+1}, ...,x_{n+1}) 
        & j < i
    \end{cases}
\end{equation}
Which is $C^\infty$ since $x_i \ne 0$. Thus, $\{\varphi_i\}$ is an atlas, and the smooth structure it defines makes $\R P^n$ a smooth manifold.

\section*{Question 5}
Let $\pi : S^n \rightarrow \R P^n$ be the map that associates to a point $x \in S^n$ the line through the origin that passes through $x$:
\begin{equation}
    \pi(x) = \{\lambda x: \lambda \ne 0\}
\end{equation} 
\subsection*{Proving $\pi$ is two-to-one}
Given two points $x_1, x_2 \in S^n$, they give the same line if for each point $\lambda_1 y_1 \in \pi(x_1)$ there exists a $\lambda_2 \ne 0$ such that:
\begin{equation}
    \label{eq12}
    \lambda_2 x_2 = \lambda_1 x_1
\end{equation}
Squaring equation \ref{eq12} results in
\begin{equation}
    \lambda_1^2 x_1^2 = \lambda_2^2 x_2^2
\end{equation}
but since $x_1, x_2 \in S^n$, their squares are $x_1^2 = x_2^2 = 1$, and thus
\begin{equation}
    \label{eq14}
    \lambda_1 = \pm \lambda_2
\end{equation}
Assigning this back into \ref{eq12} and dividing by $\lambda_1$, which is not zero by definition, gives:
\begin{equation}
    x_1 = \pm x_2
\end{equation}
Meaning there are exactly two points on $S^n$ that gives each line in $\R P^n$, so the map $\pi$ is a two-to-one map.
\subsection*{Proving $\pi$ is smooth}
The function $\pi^{-1}$ gives the set of all points responding to a specific line, i.e. the two points shown in the previous subsection $\pm x$. Therefore for a point $x \in S^n$ and an open set $V \subseteq \R P^n$ such that $\pi(x) \in V$, the set $\pi^{-1}(V)$ is an open set made of two neighborhoods around the points $\pm x$.
\newline
Now, for a point $x \in S^n$, let $(U_i, \varphi_i)$ be the chart over $S^n$ that removes the $i$-th coordinate:
\begin{equation}
    \varphi_i : S^n \rightarrow \R^n \qquad \varphi_i(x) = (x_1, ..., x_{i-1}, x_{i+1}, ..., x_n)
\end{equation}
such that $x \in U_i$ and $(V_j, \psi_j)$ a chart in the atlas defined in Question 4 over $\R P^n$ such that $\pi(x) \in V_j$. The composite function $\psi_i \circ \pi \circ \varphi_i^{-1}$ of a point $y \in \R^n$ with coordinates $y_i : i = 1,...,n$ is:
\begin{equation}
    \psi_i \pi \varphi_i^{-1} (y) =
    \left[1 - \sum_{j\ne i} y_j^2\right]^{-1/2}\left(y_1, ..., y_{i-1}, y_{i+1}, ..., y_n\right)
\end{equation}
Which is smooth for all $y \in \pi^{-1}(V_i)\cap U_i$, since the denominator is non zero. Since $\varphi_i \varphi_j^{-1}$ is smooth for ever $i,j$ as well as $\psi^i \psi_j^{-1}$, we can conclude that $\psi_i \pi \varphi_j^{-1}$ is also smooth for every $i, j$, which means it holds for all compatible atlases, and the function $\pi$ is smooth by definition.
\newline
\subsection*{proving $\pi$ is a local diffeomorphism}
The inverse of $\pi$ is locally smooth, since for every point $x \in \R P^n$ and an open set $V \subseteq \R P^n$ such that $x \in V$, $\pi(V)$ gives the set of all lines passing through the origin and through the points in $V$, which is an open set, and while $\pi^{-1}$ is not invertible over all of $S^n$, locally for each point $x \in S^n$ there is a neighborhood $V$ of $x$ such that $\pi^{-1} : \pi(V) \rightarrow V$ is one-to-one, making the composite function in equation 17 invertible with a smooth inverse, meaning that $\pi^{-1}$ is locally smooth, and according to the definition $\pi$ is locally a diffeomorphism.

\section*{Question 6}
\subsection*{The structure on $f(N)$ induced by $N$}
Let $f : N \rightarrow M$ be an embedding. Then, by definition, it is also a onto its image, and has a bijective inverse $f^{-1}$. For a chart $(U_i, \varphi_i)$ on $N$, define the chart $\left(f(U_i), \varphi_i f^{-1}\right)$ on $f(N)$, where $\varphi_i f^{-1}$ is a composition of bijections and therefore a bijection itself.
\newline
For a given atlas on $N$, the set of all such charts $\left\{\left(f(U_i), \varphi_i f^{-1}\right)\right\}$ matching charts $(U_i, \varphi_i)$ in the atlas, covers all of $f(N)$, because the atlas covers $\bigcup U_i = N$ and since $f$ is a bijection, $\bigcup f(U_i) = f\left(\bigcup U_i\right) = f(N)$.
\newline
In addition, the for each $i, j$:
\begin{equation}
    \varphi_i f^{-1} (f(U_i) \cap f(U_j)) = \varphi_i (U_i \cap U_j)
\end{equation}
And therefore, by the definition of the atlas on $N$, $\varphi_i f^{-1} (f(U_i) \cap f(U_j))$ is an open set in $\R^n$.
\newline
Finally, for every $i,j$:
\begin{equation}
    (\varphi_i f^{-1}) \circ (\varphi_j f^{-1}) = \varphi_i f^{-1} f \varphi_j^{-1} = \varphi_i \varphi_j^{-1}
\end{equation}
Which is again, by the definition of an atlas, $C^\infty$.
These three points together make $\left\{\left(f(U_i), \varphi_i f^{-1}\right)\right\}$ a smooth atlas on $f(N)$, and the structure that it generates is the smooth structure that makes $f(N)$ a manifold.
\subsection*{The structure on $f(N)$ induced by $M$}
For every point $p \in f(N)$ let $\left(U_i, \varphi_i\right)$ be a chart on $M$ such that $p \in U_i$. Let $\pi: \R^{m} \rightarrow \R^{n}$ be the projection function, where $n$ and $m$ are the dimensions of $N$ and $M$ respectively.
Thus, for an atlas $\left\{\left(U_i, \varphi_i\right)\right\}$ over $M$ the set of charts $\left\{\left(U_i \cap f(N), \pi \circ \varphi_i\right)\right\}$ covers $f(N)$, and from the rank theorem, there exist a chart $\left(V_i, \psi_i\right)$ for point $x \in N$ such that $\varphi_i \circ f \circ \psi_i^{-1} = \iota$, where $\iota$ is the inclusion map from $\R^n$ to $\R^m$ that leaves all other coordinates as zero. Acting on boths sides with the projection $\pi$, we get:
\begin{equation}
\begin{gathered}
    \pi \circ \varphi_i \circ f \circ \psi_i^{-1} = Id_{\psi_i(V_i)}
    \\
    \pi \circ \varphi_i = \psi_i \circ f^{-1}
\end{gathered}
\end{equation}
For every two charts $\left(U_i, \varphi_i\right)$ and $\left(U_j, \varphi_j\right)$ of $M$, the two induced charts on $f(N)$ are $\left(U_i \cap f(N), \pi\circ\varphi_i\right)$ and $\left(U_j \cap f(N), \pi\circ\varphi_j\right)$:
\begin{equation}
\begin{gathered}
    \pi_{m\rightarrow n}\circ\varphi_i \left(\left(U_i \cap f(N)\right)\cup\left(U_j \cap f(N)\right)\right) = 
    \\
    = \psi_i \circ f^{-1} \left(\left(U_i \cap f(N)\right)\cup\left(U_j \cap f(N)\right)\right) =
    \\
    = \psi_i  \left(f^{-1}\left(U_i \cup U_j\right)\cap N\right) = 
    \\
    = \psi_i f^{-1}\left(U_i \cup U_j\right) = 
    \\
    = \pi_{m\rightarrow n}\circ\varphi_i \left(U_i \cup U_j\right) 
\end{gathered}
\end{equation}
Since $\varphi(U_i \cup U_j)$ is an open set, and the projection is surjective and therefore is an open map, $\pi_{m\rightarrow n}\circ\varphi_i \left(U_i \cup U_j\right)$ is also an open set.
Finally, for each two charts $\left(U_i \cap f(N), \pi \circ \varphi_i\right)$ and $\left(U_j \cap f(N), \pi \circ \varphi_j\right)$, their transition map is smooth:
\begin{equation}
    \left(\pi \circ \varphi_i\right) \circ \left(\pi \circ \varphi_j\right)^{-1} = 
    \pi \circ \varphi_i \circ \varphi_j^{-1} \circ \iota
\end{equation}
Putting this all together, every smooth atlas on $M$ also induces a smooth atlas on $f(N)$, therefore the structure of $M$ induces a structure on $f(N)$, and $f(N)$ is a submanifold of $M$.

\subsection*{Proving both structures coincide}
Let $\left\{\left(f(U_i), \varphi_i f^{-1}\right)\right\}$ be an atlas induced from $N$, and let $\left\{\left(V_j, \psi_j\right)\right\}$ be an atlas induced by $M$. Since $f$ is an embedding, it is also a smooth map, which by definition means $\varphi_i \circ f^{-1} \circ \psi_j|_{f(N)}$ is also smooth. Thus, the two atlases are compatible and their structures coincide.

\section*{Question 7}
Let $\Delta = \left\{(p,p)\in M \times M : p \in M\right\}$ where $M$ is an m-dimensional smooth manifold. From question 1 it is clear that $M \times M$ is a manifold of dimension $2m$. So, to show that $\Delta$ is a submanifold of $M \times M$ we must find a smooth structure over $\Delta$ for which the inclusion map $\Delta \xhookrightarrow {} M \times M$ is an embedding.
\newline
First let us build the atlas for $\Delta$. Let $\left(U_i, \varphi_i\right)$ be a chart of $M$. Now, let $f: \Delta \rightarrow M$ be defined such that $f((p,p)) = p$. It is clear that the set of charts $\left\{\left(f^{-1}(U_i), \varphi_i\circ f\right)\right\}$ is an atlas:
\begin{itemize}
    \item It covers all of $\Delta$
    \item For every two charts $\left(f^{-1}(U_i), \varphi_i\circ f\right)$ and $\left(f^{-1}(U_i), \varphi_i\circ f\right)$, the set $\varphi_i\circ f \left(f^{-1}(U_i) \cap f^{-1}(U_j)\right) = \varphi_i \left(U_i \cap U_j\right)$ is an open set.
    \item For every two charts $\left(f^{-1}(U_i), \varphi_i\circ f\right)$ and $\left(f^{-1}(U_i), \varphi_i\circ f\right)$, their transition map $\varphi_i \circ f \circ \left(\varphi_j \circ f\right)^{-1} = \varphi_i \circ f \circ f^{-1} \circ \varphi_j^{-1} = \varphi_i \circ \varphi_j^{-1}$ is smooth.
\end{itemize}
This means $\Delta$ is a manifold with the structure generated by this atlas. It is also clear from the homeomorphism that it has the same dimension as M.
\newline
Next, let us consider the tangent space to $\Delta$. For every point $p \in M$ let $T_pM$ be the tangent space to $M$ at $p$, and let $v \in T_pM$. For every function $g \in C^{\infty}(\Delta)$, the differential of $f$, $df_p: T_p M \rightarrow T_{(p,p)\Delta}$ must follow the rule $df_p(v) = v(g \circ f) = v(g,g)$. This gives $df_p(v) = (v,v)$, and the tangent space to $\Delta$ at point $(p,p)$ is
\begin{equation}
    T_{(p,p)}\Delta = \left\{(v,v) : v \in T_pM\right\}
\end{equation}
Now, showing that the inclusion map $\theta: \Delta \rightarrow M \times M$ is an embedding is fairly simple. It is obvious that the inclusion map is smooth, and that it is a homeomorphism onto its image. So, all that is left to show is that it's an immersion. To show this, let us look at its differential $d\theta: T_{(p,p)\Delta} \rightarrow T_{(p,p)}(M \times M)$. According to proposition 3.9 in Lee's book, the differential of an inclusion map from a subset of a manifold to the manifold itself is an isomorphism, which is therefore injective. By the definition, $\theta$ is an immersion and therefore and embedding, and $\Delta$ is thus an embedded submanifold.
\newline
The map $\iota: M \rightarrow M \times M, \iota(p) = (p,p)$ is then exactly the inverse of $f$, and as we saw its differential is an isomorphism and therefore injective. It is obviously smooth and a homeomorphism onto its image, which is $\Delta$. This means that $\iota$ is an embedding.

\section*{Question 8}
Let $f : M \rightarrow N$ be a smooth immersion that is also injective, and let $M$ be compact. Since N is a manifold, and thus by definition a Hausdorff space, and since an injective function from a compact space to a Hausdorff space is a homeomorphism onto its image, $f$ is a homeomorphism onto its image and therefore an embedding.

\end{document}